\documentclass[tikz,border=15pt]{standalone}
\usepackage{ctex}
\usepackage{amsmath}
\usepackage{amssymb}
\usetikzlibrary{shapes.geometric, arrows.meta, positioning, fit, backgrounds, shadows, calc}

\pgfdeclarelayer{background}
\pgfsetlayers{background,main}

\begin{document}

% ================= [ ] 内部严禁空行，确保完美编译 =================
\begin{tikzpicture}[
    >=Stealth,
    box/.style={rectangle, rounded corners=2pt, draw=black!70, fill=white, thick, align=center, drop shadow={opacity=0.08, shadow xshift=1.5pt, shadow yshift=-1.5pt}},
    pipebox/.style={box, text width=4.8cm, minimum height=1.1cm, inner sep=6pt, font=\small},
    digibox/.style={box, draw=blue!70!black, text width=5.2cm, align=left, inner sep=8pt, font=\small},
    anabox/.style={box, draw=red!70!black, text width=5.2cm, align=left, inner sep=8pt, font=\small},
    multbox/.style={box, draw=teal!70!black, text width=2.8cm, minimum height=1.2cm, inner sep=6pt, font=\small},
    diam/.style={diamond, aspect=2.4, draw=orange!80!black, fill=orange!5, thick, text centered, inner sep=0pt, text width=1.8cm, font=\footnotesize, drop shadow={opacity=0.08}},
    % 修复：将圆角从12pt降为8pt，防止超限导致渲染穿模
    startstop/.style={rectangle, rounded corners=8pt, draw=black!80, fill=gray!15, thick, text width=2.0cm, text centered, minimum height=0.7cm, font=\small\bfseries, drop shadow={opacity=0.08}},
    phasebg/.style={rectangle, draw=#1, dashed, thick, inner sep=0pt, fill=#1!3, rounded corners=5pt}
]

    % ================= 左侧：预处理管线 (大幅拉开垂直间距) =================
    \node (start) [startstop] at (0, 0) {开始};
    
    \node (input) [pipebox] at (0, -1.8) {\textbf{解耦重构信息输入}\\[0.1cm] \footnotesize 信道状态 (CSI) \& 硬件失真参数};
    
    \node (model) [pipebox] at (0, -3.6) {\textbf{构建跨域鲁棒优化模型}\\[0.1cm] \footnotesize 融合通感双目标与多重物理约束};
    
    \node (bussgang) [pipebox] at (0, -5.4) {\textbf{统计等效线性化处理}\\[0.1cm] \footnotesize 依托 Bussgang 定理处理失真};
    
    \node (lagrange) [pipebox] at (0, -7.2) {\textbf{构建增广拉格朗日函数 $\mathcal{L}_\rho$}\\[0.1cm] \footnotesize 引入辅助变量，彻底剥离深层耦合};

    % 左侧连线
    \draw [->, line width=1.2pt, black!70] (start) -- (input);
    \draw [->, line width=1.2pt, black!70] (input) -- (model);
    \draw [->, line width=1.2pt, black!70] (model) -- (bussgang);
    \draw [->, line width=1.2pt, black!70] (bussgang) -- (lagrange);


    % ================= 右侧：ADMM 核心框架 (彻底解决重叠) =================
    \node (digital) [digibox] at (7.5, -1.2) {
        \textbf{子块 1: 常规变量交替优化}\\[0.15cm]
        $\bullet$ 更新\textbf{数字预编码}及波形辅助变量\\
        $\bullet$ 基于增广拉格朗日函数剥离耦合\\[0.1cm]
        % 安全的内部横线绘制方式
        {\color{blue!20}\rule{4.8cm}{0.8pt}}\\[0.1cm]
        {\footnotesize\color{blue!70!black} $\blacktriangleright$ \textbf{当前状态:} 固定模拟移相器变量}
    };

    % 将子块2的Y坐标下移至 -5.6，预留出完美的接驳空间
    \node (analog) [anabox] at (7.5, -5.6) {
        \textbf{子块 2: 模拟移相器优化}\\[0.15cm]
        \color{red!70!black}\textbf{突破非凸恒模约束性能瓶颈：}\\
        \color{black}$\bullet$ 等效映射至\textbf{复圆流形空间}\\
        $\bullet$ 利用 \textbf{RMCG 算法} 精确求解\\[0.1cm]
        % 安全的内部横线绘制方式
        {\color{red!20}\rule{4.8cm}{0.8pt}}\\[0.1cm]
        {\footnotesize\color{red!80!black} $\blacktriangleright$ \textbf{当前状态:} 固定常规数字变量}
    };

    \node (multi) [multbox] at (13.5, -3.4) {
        \textbf{子块 3}\\[0.15cm]
        更新拉格朗日\\对偶乘子
    };

    \node (check) [diam] at (13.5, -7.2) {残差满足\\收敛要求?};

    \node (end) [startstop] at (10.0, -9.0) {结束};


    % ================= 完美正交连线与防碰撞路由 =================
    
    % 1. 初始化变量 (降低了拐点坐标，防止横穿左侧最后一个框)
    \draw [->, line width=1.2pt, black!70, rounded corners=4pt] 
        (lagrange.east) -- (3.5, -7.2) |- (digital.west);
    \node [anchor=south, font=\footnotesize, text=black!80] at (3.5, -7.1) {初始化变量};

    % 2. 传递波形 (有了充足的空间，文字现在稳稳地待在线的右侧)
    \draw [->, line width=1.2pt, blue!70!black] 
        (digital.south) -- node[right=2pt, font=\footnotesize, text=black] {最新波形 $\mathcal{X}^{\ell+1}$} (analog.north);

    % 3. 传递流形变量 (平滑折线)
    \draw [->, line width=1.2pt, red!70!black, rounded corners=4pt] 
        (analog.east) -- (11.5, -5.6) |- node[below right=1pt, font=\footnotesize, text=black, pos=0.3] {流形变量 $\mathbf{\Theta}^{\ell+1}$} (multi.west);

    % 4. 传递最新乘子
    \draw [->, line width=1.2pt, teal!70!black, rounded corners=4pt] 
        (multi.north) |- node[above=2pt, font=\footnotesize, text=black, pos=0.75] {传递最新乘子} (digital.east);

    % 5. 乘子至判断框
    \draw [->, line width=1.2pt, dashed] (multi.south) -- (check.north);

    % 6. 迭代反馈大环线 (抬高了过顶高度 Y=0.8，绝不压迫内容)
    \draw [->, line width=1.2pt, dashed, orange!90!black, rounded corners=6pt] 
        (check.east) -- (15.5, -7.2) -- (15.5, 0.8) -- node[above=2pt, font=\footnotesize, text=black] {否 (进入 $\ell+1$ 次迭代)} (7.5, 0.8) -- (digital.north);

    % 7. 判断至结束
    \draw [->, line width=1.2pt, rounded corners=4pt] 
        (check.south) |- node[above right=2pt, font=\footnotesize] {是 (完成优化)} (end.east);


    % ================= 高度一致的宽敞背景大框 =================
    \begin{pgfonlayer}{background}
        \def\TopY{1.6}
        \def\BottomY{-10.0}
        
        % 左侧大框
        \coordinate (P1_TL) at (-2.8, \TopY);
        \coordinate (P1_BR) at (2.8, \BottomY);
        \node (phase1_bg) [phasebg=gray, fit=(P1_TL) (P1_BR)] {};
        
        % 右侧大框
        \coordinate (P2_TL) at (4.2, \TopY);
        \coordinate (P2_BR) at (16.2, \BottomY);
        \node (phase2_bg) [phasebg=blue, fit=(P2_TL) (P2_BR)] {};
    \end{pgfonlayer}

    % ================= 放置内部顶栏标题 =================
    \node [anchor=south, font=\bfseries\normalsize, text=black!70] at (0, 1.8) {跨域建模与解耦预处理};
    \node [anchor=south, font=\bfseries\normalsize, text=blue!80!black] at (10.2, 1.8) {基于 ADMM 的通感双目标流形交替优化框架};

\end{tikzpicture}

\end{document}